% GENERATED FILE. Edit canonical lessons, then run npm run generate:books.
% canonical-source-hash: fnv1a64:b480d74a627d27a3
% canonical-lessons: KA-C392-cunavane, KA-C392-mata, KA-C392-bahumata, KA-C392-pracara-madu, KA-C392-sambodhisu

\chapter{Elections}
\label{ch:ka-elections}
\hlchaptermodality{392}

\begin{chapteropening}
\textbf{By the end of this chapter:} \emph{I can say an election, a vote, a majority, to campaign and to address a crowd.}

Complete the last lesson of chapter 392: I can say an election, a vote, a majority, to campaign and to address a crowd.
\end{chapteropening}

\section[\texorpdfstring{\kn{ಚುನಾವಣೆ} (cunāvaṇe)}{ಚುನಾವಣೆ (cunāvaṇe)}]{\kn{ಚುನಾವಣೆ} (cunāvaṇe) — an election}

\label{lesson:KA-C392-cunavane}



\begin{quote}
Before the new one: say the Kannada for empowerment, then the Kannada for commitment.
\end{quote}

\subsection*{You'll want to know: \kn{ಚುನಾವಣೆ}}
\textbf{\kn{ಚುನಾವಣೆ}} — \emph{cunāvaṇe} — \textquotedblleft{}an election\textquotedblright{}.

\begin{quote}
\kn{ಮುಂದಿನ} \kn{ತಿಂಗಳು} \kn{ನಮ್ಮ} \kn{ಊರಿನಲ್ಲಿ} \kn{ಚುನಾವಣೆ} \kn{ಇದೆ}.
\end{quote}

\begin{quote}
\emph{mundina tiṅgaḷu namma ūrinalli cunāvaṇe ide.}
\end{quote}

\begin{quote}
There is an election in our town next month.
\end{quote}

\kn{ಚುನಾವಣೆಗೆ} \kn{ನಿಲ್ಲು}, to stand for election.

\begin{cousinweb}
\textbf{\kn{ಚುನಾವಣೆ}} (\emph{cunāvaṇe}) is a loan from Hindustani \emph{cunāv}, \textquotedblleft{}a choosing, an election\textquotedblright{}, from \emph{cunnā}, \textquotedblleft{}to pick\textquotedblright{}, with a Kannada noun ending.
\end{cousinweb}

\begin{grammarlens}[title={What you've built}]
One more word for giving reasons and opinions.
\end{grammarlens}

\subsection*{Your turn}
\begin{itemize}
\raggedright
  \item \emph{Say it:} \emph{cunāvaṇe}
  \item \emph{Say it:} \emph{cunāvaṇe}, once more
  \item \emph{Say it:} \emph{mundina tiṅgaḷu namma ūrinalli cunāvaṇe ide.} — There is an election in our town next month.
  \item \emph{Recall:} say the Kannada for empowerment, then the Kannada for commitment, then say \emph{cunāvaṇe} again
\end{itemize}

\begin{tcolorbox}[breakable,colback=teal!4,colframe=teal!35!black,title={Before you move on}]
What does \kn{ಚುನಾವಣೆ} mean? (\textquotedblleft{}An election\textquotedblright{}.) Say it once more.
\end{tcolorbox}
\section[\texorpdfstring{\kn{ಮತ} (mata)}{ಮತ (mata)}]{\kn{ಮತ} (mata) — a vote}

\label{lesson:KA-C392-mata}



\begin{quote}
Before the new one: say the Kannada for commitment, then the Kannada for an election.
\end{quote}

\subsection*{You'll want to know: \kn{ಮತ}}
\textbf{\kn{ಮತ}} — \emph{mata} — \textquotedblleft{}a vote\textquotedblright{}.

\begin{quote}
\kn{ಹದಿನೆಂಟು} \kn{ವರ್ಷ} \kn{ತುಂಬಿದವರಿಗೆ} \kn{ಮತ} \kn{ಹಾಕುವ} \kn{ಹಕ್ಕು} \kn{ಇದೆ}.
\end{quote}

\begin{quote}
\emph{hadineṇṭu varṣa tumbidavarige mata hākuva hakku ide.}
\end{quote}

\begin{quote}
Those who have turned eighteen have the right to vote.
\end{quote}

\kn{ಮತ} \kn{ಹಾಕು}, to cast a vote. It also means an opinion or a creed.

\begin{cousinweb}
\textbf{\kn{ಮತ}} (\emph{mata}) is borrowed from Sanskrit \emph{mata}, \textquotedblleft{}thought, opinion\textquotedblright{} (Monier-Williams), from \emph{man}, \textquotedblleft{}to think\textquotedblright{}.
\end{cousinweb}

\begin{grammarlens}[title={What you've built}]
One more word for giving reasons and opinions.
\end{grammarlens}

\subsection*{Your turn}
\begin{itemize}
\raggedright
  \item \emph{Say it:} \emph{mata}
  \item \emph{Say it:} \emph{mata}, once more
  \item \emph{Say it:} \emph{hadineṇṭu varṣa tumbidavarige mata hākuva hakku ide.} — Those who have turned eighteen have the right to vote.
  \item \emph{Recall:} say the Kannada for commitment, then the Kannada for an election, then say \emph{mata} again
\end{itemize}

\begin{tcolorbox}[breakable,colback=teal!4,colframe=teal!35!black,title={Before you move on}]
What does \kn{ಮತ} mean? (\textquotedblleft{}A vote\textquotedblright{}.) Say it once more.
\end{tcolorbox}
\section[\texorpdfstring{\kn{ಬಹುಮತ} (bahumata)}{ಬಹುಮತ (bahumata)}]{\kn{ಬಹುಮತ} (bahumata) — a majority}

\label{lesson:KA-C392-bahumata}



\begin{quote}
Before the new one: say the Kannada for an election, then the Kannada for a vote.
\end{quote}

\subsection*{You'll want to know: \kn{ಬಹುಮತ}}
\textbf{\kn{ಬಹುಮತ}} — \emph{bahumata} — \textquotedblleft{}a majority\textquotedblright{}.

\begin{quote}
\kn{ಚುನಾವಣೆಯಲ್ಲಿ} \kn{ಅವರಿಗೆ} \kn{ಬಹುಮತ} \kn{ಸಿಕ್ಕಿತು}.
\end{quote}

\begin{quote}
\emph{cunāvaṇeyalli avarige bahumata sikkitu.}
\end{quote}

\begin{quote}
They won a majority in the election.
\end{quote}

\kn{ಬಹು} (many) + \kn{ಮತ} (votes).

\begin{cousinweb}
\textbf{\kn{ಬಹುಮತ}} (\emph{bahumata}) is built from Sanskrit \emph{bahu}, \textquotedblleft{}many\textquotedblright{}, and \emph{mata}, \textquotedblleft{}opinion\textquotedblright{} (Monier-Williams); the sense \textquotedblleft{}majority\textquotedblright{} is modern.
\end{cousinweb}

\begin{grammarlens}[title={What you've built}]
One more word for giving reasons and opinions.
\end{grammarlens}

\subsection*{Your turn}
\begin{itemize}
\raggedright
  \item \emph{Say it:} \emph{bahumata}
  \item \emph{Say it:} \emph{bahumata}, once more
  \item \emph{Say it:} \emph{cunāvaṇeyalli avarige bahumata sikkitu.} — They won a majority in the election.
  \item \emph{Recall:} say the Kannada for an election, then the Kannada for a vote, then say \emph{bahumata} again
\end{itemize}

\begin{tcolorbox}[breakable,colback=teal!4,colframe=teal!35!black,title={Before you move on}]
What does \kn{ಬಹುಮತ} mean? (\textquotedblleft{}A majority\textquotedblright{}.) Say it once more.
\end{tcolorbox}
\section[\texorpdfstring{\kn{ಪ್ರಚಾರ} \kn{ಮಾಡು} (pracāra māḍu)}{ಪ್ರಚಾರ ಮಾಡು (pracāra māḍu)}]{\kn{ಪ್ರಚಾರ} \kn{ಮಾಡು} (pracāra māḍu) — to campaign; to publicise}

\label{lesson:KA-C392-pracara-madu}



\begin{quote}
Before the new one: say the Kannada for a vote, then the Kannada for a majority.
\end{quote}

\subsection*{You'll want to know: \kn{ಪ್ರಚಾರ} \kn{ಮಾಡು}}
\textbf{\kn{ಪ್ರಚಾರ} \kn{ಮಾಡು}} — \emph{pracāra māḍu} — \textquotedblleft{}to campaign; to publicise\textquotedblright{}.

\begin{quote}
\kn{ಅವಳು} \kn{ಮನೆ} \kn{ಮನೆಗೆ} \kn{ಹೋಗಿ} \kn{ಪ್ರಚಾರ} \kn{ಮಾಡಿದಳು}.
\end{quote}

\begin{quote}
\emph{avaḷu mane manege hōgi pracāra māḍidaḷu.}
\end{quote}

\begin{quote}
She campaigned from door to door.
\end{quote}

Past (I …): \kn{ಪ್ರಚಾರ} \kn{ಮಾಡಿದೆ}, \emph{pracāra māḍide}.

\kn{ಪ್ರಚಾರ} alone is publicity or a campaign.

\begin{cousinweb}
\textbf{\kn{ಪ್ರಚಾರ} \kn{ಮಾಡು}} (\emph{pracāra māḍu}) is built from Sanskrit \emph{pracāra}, \textquotedblleft{}coming forth, showing oneself\textquotedblright{} (Monier-Williams), with \emph{māḍu}, \textquotedblleft{}to do\textquotedblright{}.
\end{cousinweb}

\begin{grammarlens}[title={What you've built}]
One more word for giving reasons and opinions.
\end{grammarlens}

\subsection*{Your turn}
\begin{itemize}
\raggedright
  \item \emph{Say it:} \emph{pracāra māḍu}
  \item \emph{Say it:} \emph{pracāra māḍu}, once more
  \item \emph{Say it:} \emph{avaḷu mane manege hōgi pracāra māḍidaḷu.} — She campaigned from door to door.
  \item \emph{Recall:} say the Kannada for a vote, then the Kannada for a majority, then say \emph{pracāra māḍu} again
\end{itemize}

\begin{tcolorbox}[breakable,colback=teal!4,colframe=teal!35!black,title={Before you move on}]
What does \kn{ಪ್ರಚಾರ} \kn{ಮಾಡು} mean? (\textquotedblleft{}To campaign; to publicise\textquotedblright{}.) Say it once more.
\end{tcolorbox}
\section[\texorpdfstring{\kn{ಸಂಬೋಧಿಸು} (sambōdhisu)}{ಸಂಬೋಧಿಸು (sambōdhisu)}]{\kn{ಸಂಬೋಧಿಸು} (sambōdhisu) — to address a crowd}

\label{lesson:KA-C392-sambodhisu}



\begin{quote}
Before the new one: say the Kannada for a majority, then the Kannada for to campaign.
\end{quote}

\subsection*{You'll want to know: \kn{ಸಂಬೋಧಿಸು}}
\textbf{\kn{ಸಂಬೋಧಿಸು}} — \emph{sambōdhisu} — \textquotedblleft{}to address a crowd\textquotedblright{}.

\begin{quote}
\kn{ಚುನಾವಣೆಗೆ} \kn{ಮೊದಲು} \kn{ಅವಳು} \kn{ಊರಿನ} \kn{ಜನರನ್ನು} \kn{ಸಂಬೋಧಿಸಿದಳು}.
\end{quote}

\begin{quote}
\emph{cunāvaṇege modalu avaḷu ūrina janarannu sambōdhisidaḷu.}
\end{quote}

\begin{quote}
Before the election she addressed the townspeople.
\end{quote}

Past (I …): \kn{ಸಂಬೋಧಿಸಿದೆ}, \emph{sambōdhiside}.

The audience takes -\kn{ಅನ್ನು}. \kn{ಸಂಬೋಧನೆ}, an address.

\begin{cousinweb}
\textbf{\kn{ಸಂಬೋಧಿಸು}} (\emph{sambōdhisu}) is built on Sanskrit \emph{saṃbodhana}, \textquotedblleft{}calling to, addressing\textquotedblright{} (Monier-Williams), with the Kannada verb ending -\emph{isu}.
\end{cousinweb}

\begin{grammarlens}[title={What you've built}]
One more word for giving reasons and opinions.
\end{grammarlens}

\subsection*{Your turn}
\begin{itemize}
\raggedright
  \item \emph{Say it:} \emph{sambōdhisu}
  \item \emph{Say it:} \emph{sambōdhisu}, once more
  \item \emph{Say it:} \emph{cunāvaṇege modalu avaḷu ūrina janarannu sambōdhisidaḷu.} — Before the election she addressed the townspeople.
  \item \emph{Recall:} say the Kannada for a majority, then the Kannada for to campaign, then say \emph{sambōdhisu} again
\end{itemize}

\begin{tcolorbox}[breakable,colback=teal!4,colframe=teal!35!black,title={Before you move on}]
What does \kn{ಸಂಬೋಧಿಸು} mean? (\textquotedblleft{}To address a crowd\textquotedblright{}.) Say it once more.
\end{tcolorbox}
